\chapter{Lawvere Theories for Dataplane Components}\label{ch:lawvere}

A component is a molecule governed by equations.
A last-in first-out buffer obeys cancellation laws relating
$\mathit{push}$ and $\mathit{pop}$; a first-in first-out ring of
capacity $2^{k}$ obeys an occupancy bound on a pair of monotone
pointers; the parser obeys the laws of sequence and choice; the
transport obeys the laws of batching. This chapter collects the operations
and equations of each component into a \emph{Lawvere theory}~\cite{lawvere,
arxiv-lawvere} (\S\ref{sec:prelim-lawvere}), with three consequences. First,
the equations become theorems: an identity of the theory holds in
\emph{every} model, so a law verified once holds for every realization of
the component in Mol. Second, the theory is the instrument of
\emph{completeness}: for the ring theory the free model decides every term
identity, so the rewrite engine of \autoref{ch:rewrite} is complete on the
ring fragment. Third, the equations, read from left to right, are
\emph{rewrite rules}: they license the optimizations of a dataplane:
ring cancellation (\S\ref{sec:lawvere-ring}), parser left factoring
(\S\ref{sec:lawvere-parser}), and batch regrouping
(\S\ref{sec:lawvere-transport}): and each rule is sound because the
equation is an identity of the behavioral quotient Mol/$\approx$
(\autoref{thm:6}). The theories, their models in Mol, and the rules they
license are developed in \S\ref{sec:lawvere-ring}--\S\ref{sec:lawvere-rewrite},
with \autoref{thm:9}--\autoref{thm:12} stated and proved in place.

\paragraph{Composition convention.}
The sequential composite of molecules is written in pipeline order: $f;g$
denotes ``perform $f$, then $g$'', the composite $g \circ f$ of
\autoref{sec:prelim-cat} read in the opposite direction; the catalog
statements of \autoref{ch:intro} write this composite of the unary ring
operations with $\circ$. The tensor $f \otimes g$ is the parallel composite
of Mol (\autoref{ch:mol}), a symmetric monoidal category (\autoref{thm:3}).
Equations whose two sides realize genuinely different state spaces are
asserted in the behavioral quotient Mol/$\approx$ (\autoref{thm:6}); where
the state spaces are in bijection, they hold in Mol itself.

\section{The Stack Theory and the FIFO Ring}\label{sec:lawvere-ring}

Two buffer disciplines appear in a dataplane, and they obey different
equations. A \emph{stack} (last-in first-out) inserts and removes at the
same end: $\mathit{push}$ followed by $\mathit{pop}$ is the identity on
the element and the remaining content. A \emph{FIFO ring} of capacity
$N = 2^{k}$ inserts at a write pointer $w$ and removes at a read
pointer $r$: $\mathit{push}$ followed by $\mathit{pop}$ returns the
oldest stored element, which equals the element just pushed only when
the ring was empty. The rewrite engine of \autoref{ch:rewrite} decides
the stack fragment and the pointer fragment; occupancy of the FIFO ring
is \autoref{thm:48}.

\begin{definition}[Stack theory]
\label{def:ring-theory}
The \emph{stack theory} $\mathcal{R}$ (the LIFO buffer theory) is the
one-sorted Lawvere theory over the unary signature
$\{\mathit{push}, \mathit{pop}, \mathit{peek}\}$ presented by the equations
\begin{align}
\mathit{push};\mathit{pop} &= \mathit{id}, \label{eq:r1}\\
\mathit{pop};\mathit{push} &= \mathit{id} &&\text{(when the buffer is
non-empty)}, \label{eq:r2}\\
\mathit{peek};\mathit{pop} &= \mathit{pop}, \label{eq:r3}\\
\mathit{peek};\mathit{peek} &= \mathit{peek}, \label{eq:r4}
\end{align}
with $\mathit{id}$ the identity operation. A \emph{model} of $\mathcal{R}$
in Mol assigns to the generic object a type $B$ of buffer states, realizes
each operation as a molecule of the appropriate type, and satisfies the
equations as identities of Mol/$\approx$; the guarded equations hold on the
domains on which the operations are defined.
\end{definition}

The \emph{buffer-content model} realizes the theory on finite sequences.
Fix a type $E$ of elements; write $a \mathbin{::} s$ for the sequence with
front element $a$ and remaining content $s$. On the state space
$B = E^{*}$ of finite buffers the operations are realized by the
transformations
\[
\mathit{push} (a,s) = a \mathbin{::} s, \qquad
\mathit{pop} (a \mathbin{::} s) = (a,s), \qquad
\mathit{peek} (a \mathbin{::} s) = (a,\ a \mathbin{::} s),
\]
the last two defined on the non-empty buffers. This is the LIFO
discipline: $\mathit{push}$ inserts at the front, $\mathit{pop}$ removes
the front element, and the two are behavioral inverses up to the guards.
A FIFO ring of capacity $N = 2^{k}$ does \emph{not} satisfy
\eqref{eq:r1} on nonempty states: $\mathit{push}$ writes at $w$ and
increments $w$, $\mathit{pop}$ reads at $r$ and increments $r$, and on a
buffer that already holds an element the pop returns that older element.
The FIFO content model therefore lies outside $\mathcal{R}$. What a
power-of-two FIFO ring does satisfy is the occupancy invariant
$0 \leq w - r \leq 2^{k} - 1$ of \autoref{thm:48}, on the pointer pair
$(w, r)$.

\begin{remark}[The free ring model]
\label{rem:free-ring}
The free algebra of the $\{\mathit{push}, \mathit{pop}\}$ fragment on one
generator is the pointer line $\mathbb{Z}$: carrier
$\mathbb{Z}$, $\mathit{push} (n) = n+1$, $\mathit{pop} (n) = n-1$, generator
$0$. This models a single ring index, not FIFO content. Every cyclic
quotient $\mathbb{Z}/2^{k}$ is a finite model of the fragment, and
$\mathbb{Z}$ is the free model for which the theory is complete
(\autoref{thm:9}).
\end{remark}

\setcatalog{9}
\begin{theorem}[Stack theory equations (\autoref{thm:9})]\label{thm:9}
The stack theory equations $\mathit{push};\mathit{pop} = \mathit{id}$
and $\mathit{pop};\mathit{push} = \mathit{id}$ (when the buffer is
non-empty) hold in every model of $\mathcal{R}$, and the equations are
complete for the free pointer model $\mathbb{Z}$. They fail in the FIFO
content model on every nonempty buffer.
\end{theorem}

\begin{proof}
Read in pipeline order, the two equations are \eqref{eq:r1}--\eqref{eq:r2};
the equations of a presentation hold in every model, so it remains to
verify the intended model and the completeness claim.

\emph{The buffer-content model is a model.} On $B = E^{*}$,
$\mathit{push};\mathit{pop}$ sends $(a,s)$ to $a \mathbin{::} s$ and back
to $(a,s)$: the identity transformation on $E \otimes B$. On a non-empty
buffer $a \mathbin{::} s$, $\mathit{pop};\mathit{push}$ sends the state to
$(a,s)$ and back to $a \mathbin{::} s$: the identity wherever $\mathit{pop}$
is defined. The peek equations are verified identically:
$\mathit{peek};\mathit{pop}$ and $\mathit{pop}$ both map $a \mathbin{::} s$
to $(a,s)$, and $\mathit{peek};\mathit{peek} = \mathit{peek}$ since
$\mathit{peek}$ leaves the state unchanged.

\emph{Completeness for the free ring model.} Orient the two equations as
the rewrite rules $\mathit{push};\mathit{pop} \rightsquigarrow \mathit{id}$
and $\mathit{pop};\mathit{push} \rightsquigarrow \mathit{id}$, and consider
terms in one variable $x$. Each step deletes one $\mathit{push}$ and one
$\mathit{pop}$, so the \emph{net exponent} $e (t) = \#\mathit{push} -
\#\mathit{pop}$ is invariant, and a term with no adjacent pair of distinct
operations is a pure run $\mathit{push}^{k} (x)$ or $\mathit{pop}^{k} (x)$,
of which exactly one exists per value of $e$. Every term reduces to the
unique normal form with its net exponent, so two terms are provably equal
exactly when their net exponents agree. In the free ring model
$\mathbb{Z}$ of \autoref{rem:free-ring}, a term of net exponent $e$
evaluates to $e$, so two terms are equal in the free ring model exactly
when their net exponents agree. Provable equality and validity therefore
coincide: the equations are complete for the free pointer model, and the same
argument decides the word problem of the fragment by normal forms, the
content of the buffer rewrite system of \autoref{ch:rewrite}. In the FIFO
content model on $E^{*}$, $\mathit{push}$ appends at the tail and
$\mathit{pop}$ removes at the head, so
$\mathit{push};\mathit{pop}$ sends $(a, s)$ to
$\mathit{pop} (s \mathbin{+\!\!+} [a])$, which equals $(a, s)$ if and
only if $s$ is empty. On a nonempty buffer $s = e \mathbin{::} t$ the
composite returns $(e, t \mathbin{+\!\!+} [a])$. So \eqref{eq:r1} fails
for FIFO content, as claimed.
\end{proof}

\section{The Parser Theory}\label{sec:lawvere-parser}

Parsers in a dataplane must be deterministic: at every point the machine
makes exactly one choice, resolved by the next unconsumed token, the LL(1)
discipline. Under this discipline the combinator structure of parsing is a
Lawvere theory.

\begin{definition}[Parser theory]
\label{def:parser-theory}
The \emph{parser theory} $\mathcal{P}$ is the Lawvere theory over the
binary operations $\mathit{seq}$ (sequence) and $\mathit{choice}$, and the
nullary operations $\varepsilon$ (the empty parser) and $\mathit{fail}$
(the failing parser), presented by the equations
\begin{align}
\mathit{seq} (\mathit{seq} (p,q),r) &= \mathit{seq} (p,\mathit{seq} (q,r))
&&\text{(sequence associative)},\\
\mathit{seq} (\varepsilon,p) &= p = \mathit{seq} (p,\varepsilon)
&&\text{(identity laws)},\\
\mathit{choice} (\mathit{choice} (p,q),r) &= \mathit{choice} (p,\mathit{choice} (q,r))
&&\text{(choice associative)},\\
\mathit{choice} (p,\mathit{fail}) &= p = \mathit{choice} (\mathit{fail},p)
&&\text{(failure identity)},\\
\mathit{seq} (p,\mathit{choice} (q,r)) &= \mathit{choice} (\mathit{seq} (p,q),\mathit{seq} (p,r))
&&\text{(left distribution)}.
\end{align}
The last equation is \emph{left distribution}: sequence distributes over
choice on the left, and read from right to left it is \emph{left
factoring}, the extraction of the common prefix $p$.
\end{definition}

A \emph{parser} in Mol is a molecule $P : T^{*} \to T^{*} \otimes R$ over
a token type $T$ and a result type $R$: it consumes a prefix of the input
stream and returns the remaining stream with a result, failing (via the
effectus structure of \autoref{ch:mol}) outside its domain. Sequence is the
Mol composite $p;q$; choice runs $p$'s lookahead test and takes the
matching branch. Each choice is resolved by the next unconsumed token, so
the models of $\mathcal{P}$ are the deterministic parsers of the LL(1)
discipline, and the equations of \autoref{def:parser-theory} are identities
of this semantics (\autoref{thm:10}).

\setcatalog{10}
\begin{theorem}[Parser theory (\autoref{thm:10})]\label{thm:10}
In the parser theory, sequence distributes over choice, choice is
associative, the identity laws hold, and left factoring is sound: the
rewrite
\[
\mathit{choice} (\mathit{seq} (p,q),\ \mathit{seq} (p,r)) \rightsquigarrow
\mathit{seq} (p,\ \mathit{choice} (q,r))
\]
preserves behavior.
\end{theorem}

\begin{proof}
Work in the stream semantics of \S\ref{sec:lawvere-parser}: a parser
denotes, for every input stream, the pair (remaining stream, result), or
failure. Sequence is composition of these partial transformations and
choice is the union over the branch selected by the lookahead test, so it
suffices to compare the two sides of each equation on every stream.

Associativity of $\mathit{seq}$ and of $\mathit{choice}$ is that of the
underlying composition and of the ternary ``first, else'' decision. The
identity laws hold because $\varepsilon$ consumes nothing and always
succeeds, while $\mathit{fail}$ never succeeds: both
$\mathit{seq} (\varepsilon,p) = p = \mathit{seq} (p,\varepsilon)$ and
$\mathit{choice} (p,\mathit{fail}) = p = \mathit{choice} (\mathit{fail},p)$
are immediate. For left distribution, fix $p,q,r$. On the left, the input
stream is first processed by $p$, and the choice between $q$ and $r$ is
resolved by the first unconsumed token. On the right, the same stream is
accepted exactly when it is accepted by $\mathit{seq} (p,q)$ or
$\mathit{seq} (p,r)$, i.e.\ exactly when $p$ succeeds on it and $q$ or $r$
succeeds on the remainder; and the result is $p$'s result followed by the
branch result, the same on both sides. The two sides therefore accept the
same streams and return the same remainder and result, hence are equal in
Mol/$\approx$. Soundness of left factoring is the same equation read right
to left: the rewrite replaces a term by an equal one, so behavior is
preserved on every input.
\end{proof}

\section{The Transport Theory}\label{sec:lawvere-transport}

The transport moves messages across a boundary: a socket, a ring, a
device: and the fixed cost of each crossing makes batching the central
optimization (\autoref{ch:operad}). The transport theory makes precise the
equations that batching relies on.

\begin{definition}[Transport theory]
\label{def:transport-theory}
Fix a type $M$ of messages. The \emph{transport theory} $\mathcal{T}$ is
the Lawvere theory over the unary operations $\mathit{recv}$ and
$\mathit{send}$ and the $n$-ary operations $\mathit{batch}_{n}$
($n \geq 1$), presented by the equations
\[
\mathit{batch}_{n} (\mathit{recv},\ldots,\mathit{recv})
= \underbrace{\mathit{recv} \otimes \cdots \otimes \mathit{recv}}_{n},
\qquad
\mathit{batch}_{n} (\mathit{send},\ldots,\mathit{send})
= \underbrace{\mathit{send} \otimes \cdots \otimes \mathit{send}}_{n},
\]
and by the \emph{interchange law}
\[
(f \otimes g);(h \otimes k) = (f;h) \otimes (g;k)
\]
for all molecules $f,g,h,k$ of matching types. A \emph{model} of
$\mathcal{T}$ in Mol realizes $\mathit{recv}$ and $\mathit{send}$ as
transport molecules and $\mathit{batch}_{n}$ as the $n$-batch construction
of \autoref{ch:operad}; the batch equations state that batching is
tensoring, which holds for pure operations: the datagram discipline, in
which successive receives are independent: and the interchange law holds
in every model.
\end{definition}

\setcatalog{11}
\begin{theorem}[Transport theory (\autoref{thm:11})]\label{thm:11}
In the transport theory, $\mathit{batch\_recv} = \mathit{recv} \otimes
\cdots \otimes \mathit{recv}$, and the interchange law
$(f \otimes g);(h \otimes k) = (f;h) \otimes (g;k)$ holds.
\end{theorem}

\begin{proof}
The batch equation is the defining equation of
\autoref{def:transport-theory}, and it holds in the intended model: by the
batch construction of \autoref{ch:operad}, the $n$-batch of a molecule
applies its transition $n$ times, once per input, threading the state; a
pure molecule has a single-element state space, so there is no state to
thread, and its $n$-fold tensor applies the same transition to each
coordinate independently. Both sides denote the same transformation on the
tuple $A^{n}$, so $\mathit{batch}_{n} (\mathit{recv},\ldots,\mathit{recv})
= \mathit{recv} \otimes \cdots \otimes \mathit{recv}$, and likewise for
$\mathit{send}$.

For the interchange law, let $f : A \to B$, $h : B \to C$, $g : A' \to B'$,
$k : B' \to C'$; both sides are morphisms $A \otimes A' \to C \otimes C'$.
On an input $(a,a')$, the left side runs $f$ then $h$ on the first
component and $g$ then $k$ on the second; the right side runs $f;h$ and
$g;k$ on the two components. Both traverse the same transitions per
component, and the state spaces $(S_{f} \times S_{g}) \times (S_{h} \times
S_{k})$ and $(S_{f} \times S_{h}) \times (S_{g} \times S_{k})$ correspond
under the reassociation bijection, so the machines are equal in Mol: the
bifunctoriality of the tensor of a symmetric monoidal category
(\autoref{thm:3}, \cite{maclane}) in pipeline order.
\end{proof}

\section{Models in Mol and the Licensed Rewrite Rules}\label{sec:lawvere-rewrite}

Each theory of this chapter is realized by the corresponding molecule:
the ring buffer of \S\ref{sec:lawvere-ring}, the parser of
\S\ref{sec:lawvere-parser}, the transport of
\S\ref{sec:lawvere-transport}: and each equation, read from left to
right, is a rewrite rule available to the rewrite engine of
\autoref{ch:rewrite}. Since every equation is an identity in Mol/$\approx$,
every rule preserves behavior, and the termination and confluence results
of \autoref{ch:rewrite} (\autoref{thm:13}--\autoref{thm:19}) turn the rules into a decision
procedure. The central ring rule is ring cancellation: a $\mathit{push}$
immediately followed by a $\mathit{pop}$ is behaviorally invisible and can
be deleted.

\setcatalog{12}
\begin{theorem}[Stack cancellation (\autoref{thm:12})]\label{thm:12}
An adjacent pair $\mathit{push};\mathit{pop}$ is removable from any stack
term with behavior preserved: for every context $C$ in the theory,
$C[\mathit{push};\mathit{pop}] = C[\mathit{id}]$ holds in every model of
$\mathcal{R}$. The rule does not apply to FIFO content terms.
\end{theorem}

\begin{proof}
By \eqref{eq:r1}, $\mathit{push};\mathit{pop} = \mathit{id}$ in every
model. Contexts are built from the term constructors: composition with a
constant and tensor with a constant: and each preserves equality: if
$u = v$ then $w;u = w;v$, $u;w = v;w$, and $w \otimes u = w \otimes v$ for
every $w$. Induction on the construction of $C$ gives
$C[\mathit{push};\mathit{pop}] = C[\mathit{id}]$ in every model, and equal
morphisms of Mol/$\approx$ are behaviorally equivalent, so deleting the
pair: in particular from a pipeline in which a pushed element is popped
immediately: preserves behavior. The guarded equation \eqref{eq:r2}
yields the analogous rule: $\mathit{pop};\mathit{push}$ is removable
wherever the buffer is known to be non-empty.
\end{proof}

\begin{table}[htbp]
\centering
\small
\caption{The component theories, their laws, and the rewrite rules they
license. Every rule preserves behavior, because each law is an identity of
Mol/$\approx$.}
\label{tab:lawvere-rules}
\begin{tabular}{@{}llp{6.5cm}@{}}
\toprule
Theory & Law & Licensed rewrite rule\\
\midrule
Stack $\mathcal{R}$ & $\mathit{push};\mathit{pop} = \mathit{id}$
(\ref{thm:9}, \ref{thm:12}) & $\mathit{push};\mathit{pop} \rightsquigarrow
\mathit{id}$: LIFO cancellation.\\
Stack $\mathcal{R}$ & $\mathit{pop};\mathit{push} = \mathit{id}$, guarded
(\ref{thm:9}) & $\mathit{pop};\mathit{push} \rightsquigarrow
\mathit{id}$ wherever the stack is known non-empty.\\
FIFO ring & $0 \leq w-r \leq 2^{k}-1$ (\ref{thm:48}) & occupancy invariant;
bitmask indexing is sound exactly on this bound.\\
Parser $\mathcal{P}$ & left distribution (\ref{thm:10}) &
$\mathit{choice} (p;q,\ p;r) \rightsquigarrow p;\mathit{choice} (q,r)$:
left factoring.\\
Parser $\mathcal{P}$ & choice associative (\ref{thm:10}) &
$\mathit{choice} (\mathit{choice} (p,q),r) \rightsquigarrow
\mathit{choice} (p,\mathit{choice} (q,r))$: reassociation.\\
Transport $\mathcal{T}$ & batch $=$ tensor (\ref{thm:11}) &
$\mathit{batch}_{n} (\mathit{recv}) \rightsquigarrow
\mathit{recv} \otimes \cdots \otimes \mathit{recv}$: flattening.\\
Transport $\mathcal{T}$ & interchange (\ref{thm:11}) &
$(f \otimes g);(h \otimes k) \rightsquigarrow (f;h) \otimes (g;k)$:
independent branches.\\
\bottomrule
\end{tabular}
\end{table}

\begin{remark}
The three theories present components as algebraic structures in the sense
of functorial semantics~\cite{lawvere, arxiv-lawvere}: every implementation
of a component is a model, and the equational theory of the component is
the theory presented here. The conformance tests of the software standard
assert the laws of \autoref{tab:lawvere-rules} directly on the molecule
(\autoref{app:b}).
\end{remark}
